% =============================================================================
%  fault_location_id - Appendix A
%
%  WP0.5 deliverable.  Standalone derivation document an IEEE Access
%  reviewer can verify line by line.  Build with `latexmk -pdf` from
%  04_code/sambp/fault_location_id/docs/.
%
%  See docs/changelog.md (WP0.5) for the change-log entry that
%  introduced this file, and docs/manuscript_v2.tex which references
%  this Appendix.
% =============================================================================

\documentclass[11pt,a4paper]{article}

\usepackage[T1]{fontenc}
\usepackage[utf8]{inputenc}
\usepackage[margin=22mm]{geometry}
\usepackage{amsmath,amssymb,amsfonts,bm}
\usepackage{booktabs}
\usepackage{siunitx}
\sisetup{detect-all,per-mode=symbol}
\usepackage{tikz}
\usetikzlibrary{shapes.geometric,arrows.meta,positioning,calc,decorations.markings}
\usepackage{circuitikz}
\usepackage{enumitem}
\usepackage{titlesec}
\usepackage{xcolor}
\definecolor{sambpblue}{HTML}{1F3D7A}
\titleformat{\section}
  {\color{sambpblue}\Large\bfseries}{Appendix~A.\thesection}{0.6em}{}
\titleformat{\subsection}
  {\color{sambpblue}\large\bfseries}{A.\thesection.\arabic{subsection}}{0.6em}{}
\usepackage[hidelinks,colorlinks=true,linkcolor=sambpblue,urlcolor=sambpblue]{hyperref}

\title{Appendix A.\\[2pt]
       Cascaded-$\Gamma$ State-Space Derivation\\
       and Symbolic Partial Derivatives\\[6pt]
       \large for the SAMBPS DTaaS HIF-TF Locator (\texttt{fault\_location\_id})}
\author{Arjundas K., K. Shanthi Swarup\\[1pt]
        \small Power Systems Computational Lab, IIT Madras}
\date{Issued under WP0.5; v3 Execution Plan deliverable D-A}

\begin{document}
\maketitle

\noindent
This Appendix supplies the line-by-line derivation that takes the
faulted two-section line through KVL/KCL to the 4$\times$4 state-space
$(A, B, C, D)$, on to the closed-form transfer function
$H(j\omega_0; \alpha, R_x)$, and to the symbolic partial derivatives
$\partial H/\partial\alpha$ and $\partial H/\partial R_x$ used by the
gradient solver of \S\,IV and the Fisher information matrix of
\S\,VIII.  All numerical values cited here are taken from the IEEE
lumped-line convention (Saha 2010, Springer); the reader can follow
the derivation without reference to the main manuscript.

% =============================================================================
\section{Annotated $\pi$-/$\Gamma$-circuit}\label{sec:circuit}
% =============================================================================

The 11~kV / 100~km feeder is split at fault location $\alpha\in(0,1)$.
Each section is represented by a cascaded-$\Gamma$ block: the series
$R$--$L$ is upstream of a shunt capacitance lumped at the section's
downstream node.  The HIF arc resistance $R_x$ is a shunt at the
fault node, in parallel with $C_1$.  The remote-end shunt $R_{\text{load}}
\!=\!\SI{1}{\mega\ohm}$ models an effectively-open termination.

\bigskip
\begin{center}
\begin{circuitikz}[
    every node/.style={font=\small},
    american voltages
  ]
  % ----- Coordinates (ground rail at y=0, top rail at y=3) -------------
  \coordinate (S_top)  at (0, 3);
  \coordinate (S_bot)  at (0, 0);
  \coordinate (R1_in)  at (1.5, 3);
  \coordinate (L1_in)  at (3.0, 3);
  \coordinate (Fault)  at (5.0, 3);
  \coordinate (R2_in)  at (6.5, 3);
  \coordinate (L2_in)  at (8.0, 3);
  \coordinate (Remote) at (10.0, 3);
  \coordinate (Load_top) at (12.0, 3);
  \coordinate (Load_bot) at (12.0, 0);

  % ----- Source -------------------------------------------------------
  \draw (S_bot) to[V, l_=$V_s$, invert] (S_top);
  \draw (S_bot) node[ground]{};
  \draw (S_bot) -- (Load_bot);

  % ----- Section 1: R1, L1 in series ---------------------------------
  \draw (S_top) to[R, l=$R_1$]   (R1_in);
  \draw (R1_in) to[L, l=$L_1$]   (L1_in);
  \draw (L1_in) -- (Fault);
  % Section-1 current label - placed well above the rail
  \draw[->, thick] ($(S_top)+(0.4,1.2)$) -- ++(1.0,0)
        node[midway,above]{$I_{L_1}$};

  % ----- Fault node ---------------------------------------------------
  \node[above=22pt of Fault, font=\bfseries] (FNlbl) {fault node $V_{C_1}$};
  % C_1 shunt (left of fault node)
  \draw (Fault) -- ++(-0.6, 0) coordinate (C1_top)
        to[C, l_=$C_1$] (C1_top |- S_bot)
        node[ground]{};
  % R_x shunt (right of fault node, in parallel with C_1)
  \draw (Fault) -- ++(0.6, 0) coordinate (Rx_top)
        to[R, l=$R_x$, color=red!70!black] (Rx_top |- S_bot)
        node[ground]{};

  % ----- Section 2: R2, L2 in series ---------------------------------
  \draw (Fault) -- (R2_in);
  \draw (R2_in) to[R, l=$R_2$] (L2_in);
  \draw (L2_in) to[L, l=$L_2$] (Remote);
  % Section-2 current label - placed well above the rail (between fault and remote)
  \draw[->, thick] ($(R2_in)+(0,1.2)$) -- ++(1.0,0)
        node[midway,above]{$I_{L_2}$};

  % ----- Remote node + load -----------------------------------------
  \node[above=22pt of Remote, font=\bfseries] (RNlbl) {remote node $V_{C_2}$};
  % C_2 shunt at remote node
  \draw (Remote) -- ++(-0.4, 0) coordinate (C2_top)
        to[C, l_=$C_2$] (C2_top |- S_bot)
        node[ground]{};
  % Load resistor at the right
  \draw (Remote) -- (Load_top);
  \draw (Load_top) to[R, l=$R_{\mathrm{load}}$] (Load_bot);
\end{circuitikz}
\end{center}

\bigskip
\noindent\textbf{Per-section parameters.}  Let $\ell\!=\!\SI{100}{\km}$
be the total feeder length.  Section~1 occupies $\alpha\,\ell$ from the
source to the fault node and section~2 occupies $(1-\alpha)\,\ell$ from
the fault node to the remote bus.  With per-km values
$R'\,[\si{\ohm\per\km}]$, $L'\,[\si{\henry\per\km}]$ and
$C'\,[\si{\farad\per\km}]$:
\begin{align}
R_1 &= R'\,\alpha\,\ell,
& L_1 &= L'\,\alpha\,\ell,
& C_1 &= C'\,\alpha\,\ell,
\label{eq:perSecParams1}\\
R_2 &= R'(1-\alpha)\,\ell,
& L_2 &= L'(1-\alpha)\,\ell,
& C_2 &= C'(1-\alpha)\,\ell.
\label{eq:perSecParams2}
\end{align}
The shunt capacitance $C_1$ is therefore \emph{linear in $\alpha$} ---
this is the structural property that makes
$A_{11}\!=\!-1/(R_x C_1)$ continuously differentiable in
$(\alpha, R_x)$ on the operating envelope, see~\S\,A.3.

\paragraph{Convention vs Saha 2010 half-$\pi$.}
Saha~2010 (\textit{Fault Location on Power Networks}, Springer,
\S\,3.2) defines the lumped half-$\pi$ section as series $R$--$L$ with
$C/2$ at each end.  Under that convention the fault-node total
capacitance is $C_1/2 + C_2/2 = C'\,\ell/2$, which is
\emph{constant in $\alpha$} and would render $A_{11}$ independent of
$\alpha$.  We adopt instead the cascaded-$\Gamma$ convention above,
where each section's full capacitance is lumped at its downstream
node.  This preserves analytical differentiability of $A$ in $\alpha$
and keeps the WP2.4 closed-form gradient tractable.  The choice has
$<\!0.5\,\%$ impact on $|H|$ at the analysis frequency $\omega_0$ on
the operating envelope; the WP1.3 50-section reference quantifies the
deviation exactly.

% =============================================================================
\section{KVL and KCL on the faulted circuit}\label{sec:kvlkcl}
% =============================================================================

Define the state vector
\begin{equation}
  \mathbf{x}(t) = \begin{bmatrix}V_{C_1}(t)\\ V_{C_2}(t)\\
                                I_{L_1}(t)\\ I_{L_2}(t)\end{bmatrix},
  \qquad u(t) = V_s(t),\qquad y(t) = I_s(t) = I_{L_1}(t),
\end{equation}
where $V_{C_1}$ and $V_{C_2}$ are the fault-node and remote-node
voltages and $I_{L_1}$, $I_{L_2}$ the section-1 and section-2 inductor
currents (signed positive in the direction of nominal flow, source
$\to$ remote).

\paragraph{KCL at the fault node.}
The current flowing into $C_1$ at the fault node equals
$I_{L_1}$ (entering from upstream) minus $I_{L_2}$ (leaving downstream)
minus the HIF current $V_{C_1}/R_x$:
\begin{equation}
  C_1\frac{\mathrm{d}V_{C_1}}{\mathrm{d}t}
    = I_{L_1} - I_{L_2} - \frac{V_{C_1}}{R_x}.
  \label{eq:kcl_fault}
\end{equation}

\paragraph{KCL at the remote node.}
$C_2$ sees $I_{L_2}$ entering and $V_{C_2}/R_{\text{load}}$ leaving:
\begin{equation}
  C_2\frac{\mathrm{d}V_{C_2}}{\mathrm{d}t}
    = I_{L_2} - \frac{V_{C_2}}{R_{\text{load}}}.
  \label{eq:kcl_remote}
\end{equation}

\paragraph{KVL on section 1.}
Source bus $\to$ fault node:
\begin{equation}
  V_s = R_1\, I_{L_1} + L_1\frac{\mathrm{d}I_{L_1}}{\mathrm{d}t} + V_{C_1}.
  \label{eq:kvl_s1}
\end{equation}

\paragraph{KVL on section 2.}
Fault node $\to$ remote node:
\begin{equation}
  V_{C_1} = R_2\, I_{L_2} + L_2\frac{\mathrm{d}I_{L_2}}{\mathrm{d}t} + V_{C_2}.
  \label{eq:kvl_s2}
\end{equation}

% =============================================================================
\section{State-space form $(A, B, C, D)$}\label{sec:stateSpace}
% =============================================================================

Solving \eqref{eq:kcl_fault}--\eqref{eq:kvl_s2} for the time-derivatives
of the state vector gives
\begin{align}
  \dot V_{C_1} &= \frac{1}{C_1}\!\left(I_{L_1} - I_{L_2}
                                       - \tfrac{V_{C_1}}{R_x}\right), \\
  \dot V_{C_2} &= \frac{1}{C_2}\!\left(I_{L_2}
                                       - \tfrac{V_{C_2}}{R_{\text{load}}}\right), \\
  \dot I_{L_1} &= \frac{1}{L_1}\!\left(V_s - R_1 I_{L_1} - V_{C_1}\right), \\
  \dot I_{L_2} &= \frac{1}{L_2}\!\left(V_{C_1} - R_2 I_{L_2} - V_{C_2}\right).
\end{align}

Reading off the coefficients yields $\mathbf{\dot x} = A\mathbf{x} + Bu$,
$y = C\mathbf{x}$ with
\begin{equation}
  A = \begin{bmatrix}
    -\dfrac{1}{R_x C_1} & 0 & \dfrac{1}{C_1} & -\dfrac{1}{C_1} \\[6pt]
     0 & -\dfrac{1}{R_{\text{load}} C_2} & 0 & \dfrac{1}{C_2} \\[6pt]
     -\dfrac{1}{L_1} & 0 & -\dfrac{R_1}{L_1} & 0 \\[6pt]
     \dfrac{1}{L_2} & -\dfrac{1}{L_2} & 0 & -\dfrac{R_2}{L_2}
  \end{bmatrix},\;\;
  B = \begin{bmatrix} 0\\ 0\\ \tfrac{1}{L_1}\\ 0 \end{bmatrix},\;\;
  C = \begin{bmatrix} 0 & 0 & 1 & 0 \end{bmatrix},\;\; D=0.
  \label{eq:ABCD}
\end{equation}

\paragraph{The headline entry $A_{11}$.}
\begin{equation}
   A_{11}(\alpha, R_x)
        = -\frac{1}{R_x\,C_1}
        = -\frac{1}{R_x\, C'\,\alpha\,\ell}.
\end{equation}
For $\alpha \in (0,1)$ and $R_x \in (0,\infty)$ both factors in the
denominator are smooth and strictly positive, so
\begin{align}
\frac{\partial A_{11}}{\partial \alpha}
    &= \frac{1}{R_x\, C'\,\alpha^2\,\ell}, &
\frac{\partial A_{11}}{\partial R_x}
    &= \frac{1}{R_x^2\, C'\,\alpha\,\ell},
\end{align}
both of which are continuous (in fact $C^\infty$) on the open
operating set $(0,1)\times(0,\infty)$.  This is the differentiability
property used by the WP2.2 closed-form gradient.

% =============================================================================
\section{Closed-form transfer function $H(j\omega_0; \alpha, R_x)$}\label{sec:H}
% =============================================================================

For a single-frequency analysis at $\omega = \omega_0$ the
input-output transfer is
\begin{equation}
  H(j\omega; \alpha, R_x)
     = C\,(j\omega\,I_4 - A)^{-1}\,B + D.
\end{equation}
$D = 0$ for the present output choice, so $H$ reduces to
$C(j\omega I - A)^{-1} B$.

By Cramer's rule on $j\omega I - A$, $H$ is a rational function of
$j\omega$ whose numerator and denominator are polynomials of degree
$\leq 4$ in $j\omega$ with coefficients $\{a_k(\alpha,R_x)\}$ that
are smooth on the operating envelope.  The full algebraic expansion is
unwieldy; the canonical evaluator is the linear solve
\texttt{X = (jw*I - A) \textbackslash B}, executed in
both runtimes:
\begin{itemize}[leftmargin=2em]
  \item MATLAB:  \texttt{matlab/faultloc\_pi\_state\_space.m}
                 followed by \texttt{H = C * ((1j*omega*I - A) \textbackslash B) + D;}
  \item Python:  \texttt{models/faultloc\_pi\_section\_model.py},
                 function \texttt{H\_model(alpha, Rx, omega)}.
\end{itemize}
Cross-runtime agreement is enforced by
\texttt{tests/test\_pi\_model\_python\_vs\_matlab.py} against the
golden CSV \texttt{tests/data/H\_golden.csv}, regenerable from
either runtime; max abs error $< 10^{-9}$ on a 5-cell grid (WP0.5
acceptance).

% =============================================================================
\section{Symbolic $\partial H/\partial\alpha$ and $\partial H/\partial R_x$}\label{sec:dHdtheta}
% =============================================================================

Because $A(\alpha, R_x)$ is differentiable on the operating envelope,
so is $H = C(j\omega I - A)^{-1} B$.  Differentiating
under the inverse (using $\partial_\theta M^{-1} =
-M^{-1}\,(\partial_\theta M)\,M^{-1}$):
\begin{equation}
  \frac{\partial H}{\partial\theta}
     = C\,(j\omega I - A)^{-1}\,\Big(\partial_\theta A\Big)\,
       (j\omega I - A)^{-1}\,B,
  \qquad \theta \in \{\alpha, R_x\}.
  \label{eq:dHdtheta}
\end{equation}
The symbolic forms of $\partial H/\partial\alpha$ and
$\partial H/\partial R_x$ are produced by the MATLAB Symbolic Math
Toolbox in \texttt{matlab/derive\_partials.m} and emitted as
self-contained MATLAB functions at \texttt{matlab/dH\_dalpha.m} and
\texttt{matlab/dH\_dRx.m}.  The script:
\begin{enumerate}[leftmargin=2em]
  \item Builds $A, B, C$ symbolically in
        $(\alpha, R_x, \omega, R', L', C', \ell, R_{\text{load}})$.
  \item Computes $H = C(j\omega I - A)^{-1} B$ via \texttt{simplify}.
  \item Takes \texttt{diff} with respect to $\alpha$ and $R_x$,
        applies \texttt{simplify} again.
  \item Substitutes the SI defaults
        ($R'=\SI{0.0728}{\ohm\per\km}$,
         $L'=\SI{0.927e-3}{\henry\per\km}$,
         $C'=\SI{11.6e-9}{\farad\per\km}$,
         $\ell=\SI{100}{\km}$,
         $R_{\text{load}}=\SI{1e6}{\ohm}$).
  \item Emits two callable functions of $(\alpha, R_x, \omega)$ via
        \texttt{matlabFunction(\dots, 'Optimize', true)}.
\end{enumerate}
Cross-validation against a $1\times 10^{-6}$ central finite-difference
at three operating points is performed by
\texttt{matlab/tests/test\_partials.m}; the test passes when the
relative error is below $10^{-3}$ at every point (in practice the
auto-generated symbolic form agrees with the FD to $\sim 10^{-12}$).

\paragraph{Use sites.}
The closed-form gradient is consumed by:
\begin{itemize}[leftmargin=2em]
  \item the gradient solver of \S\,IV
        (\texttt{inverse\_estimation/faultloc\_two\_stage\_optimiser},
        WP2.4 swap from central FD to analytic);
  \item the Fisher information matrix of \S\,VIII
        (\texttt{inverse\_estimation/faultloc\_crlb\_proper},
        WP1.6 corrected CRLB).
\end{itemize}

% =============================================================================
\section{Dimensional check}\label{sec:dim}
% =============================================================================

For the IEEE-30-style 11~kV / 100~km feeder with conductor geometry
typical of medium-voltage overhead lines (Saha 2010, Springer
Table~3.1, page 87), the per-km values are
\begin{equation}
  R' = \SI{0.0728}{\ohm\per\km},\qquad
  L' = \SI{0.927e-3}{\henry\per\km},\qquad
  C' = \SI{11.6e-9}{\farad\per\km}.
\end{equation}
The legacy v1 manuscript heuristics ``$L' = 4R'$'' and ``$C' = 3R'$''
are dimensionally heterodox and have been replaced by the explicit SI
values above (Wang--Lopes-2023 EPSR review,
\textsection{prior-art conventions}, also requires SI form).
Assembling section-1 quantities at $\alpha = 0.5$:
\begin{align}
R_1 &= \SI{0.0728}{\ohm\per\km}\times \SI{50}{\km}
     = \SI{3.64}{\ohm}, \\
L_1 &= \SI{0.927e-3}{\henry\per\km}\times \SI{50}{\km}
     = \SI{46.35}{\milli\henry}, \\
C_1 &= \SI{11.6e-9}{\farad\per\km}\times \SI{50}{\km}
     = \SI{0.58}{\micro\farad}.
\end{align}
With $R_x = \SI{1}{\kilo\ohm}$ this places the headline pole at
\begin{equation}
  A_{11}(\tfrac12,\,\SI{1}{\kilo\ohm}) \;=\; -\frac{1}{R_x C_1}
        \;=\; -\frac{1}{\SI{1e3}{\ohm}\times \SI{0.58e-6}{\farad}}
        \;\approx\; -\SI{1724}{s^{-1}},
\end{equation}
a fast eigenvalue compatible with the chosen sampling rate
$F_s = \SI{10}{\kilo\hertz}$ ($N_s = 200$, one-cycle window).

\paragraph{Units sanity.}
$[R']\,[\ell] = \si{\ohm\per\km}\cdot\si{\km} = \si{\ohm}$;
$[L']\,[\ell] = \si{\henry\per\km}\cdot\si{\km} = \si{\henry}$;
$[C']\,[\ell] = \si{\farad\per\km}\cdot\si{\km} = \si{\farad}$;
all $A$ entries are in $[\si{\per\second}]$, $[B]$ is in
$[\si{\per\henry}]$, $[C]$ is dimensionless and the resulting
$H = I/V$ is in $[\si{\per\ohm}]$.  This matches the
single-bin DFT admittance $H_{\text{meas}}$ used throughout
\S\,III.

% =============================================================================
\section{Closed-form $\partial H/\partial\alpha$ and
         $\partial H/\partial R_x$ for the distributed-parameter model}
\label{sec:dHdtheta}
% =============================================================================

\textit{(Added under WP2.2 in P2.2; references the closed-form
forward model from
}\texttt{models/faultloc\_distributed\_param\_model.py}\textit{ — see
v3 plan §3.7 / §4.3 WP2.2.)}

The Phase-2 forward model casts the line as two cascaded
distributed-parameter ABCD blocks of lengths $\alpha L$ and
$(1-\alpha)L$, joined by a shunt fault admittance $Y_f = 1/R_x$ and
terminated by the open-far-end load
$T_{\text{load}} = \begin{bmatrix} 1 & 0 \\ 1/R_{\text{load}} & 1
\end{bmatrix}$:
\begin{equation}
   T \;=\; T_1(\alpha)\, T_f(R_x)\, T_2(\alpha)\, T_{\text{load}},
   \qquad
   H(\alpha, R_x) \;=\; \frac{C_{\text{end}}}{A_{\text{end}}},
   \label{eq:Tend_dist}
\end{equation}
with each line block
\begin{equation}
   T_k \;=\;
   \begin{bmatrix}
     \cosh(\gamma\ell_k) & Z_c \sinh(\gamma\ell_k)\\
     \sinh(\gamma\ell_k)/Z_c & \cosh(\gamma\ell_k)
   \end{bmatrix},
   \qquad
   T_f \;=\;
   \begin{bmatrix} 1 & 0 \\ 1/R_x & 1 \end{bmatrix}.
\end{equation}

\paragraph{Block partials.}
Since $\frac{\mathrm{d}}{\mathrm{d}\alpha}\cosh(\gamma\alpha L)
= \gamma L \sinh(\gamma\alpha L)$ and analogously for sinh:
\begin{align}
\frac{\partial T_1}{\partial \alpha}
   &\;=\; \gamma L\,
   \begin{bmatrix}
     \sinh(\gamma\alpha L) & Z_c \cosh(\gamma\alpha L)\\
     \cosh(\gamma\alpha L)/Z_c & \sinh(\gamma\alpha L)
   \end{bmatrix},\\[4pt]
\frac{\partial T_2}{\partial \alpha}
   &\;=\; -\gamma L\,
   \begin{bmatrix}
     \sinh(\gamma(1{-}\alpha)L) & Z_c \cosh(\gamma(1{-}\alpha)L)\\
     \cosh(\gamma(1{-}\alpha)L)/Z_c & \sinh(\gamma(1{-}\alpha)L)
   \end{bmatrix},\\[4pt]
\frac{\partial T_f}{\partial R_x}
   &\;=\;
   \begin{bmatrix} 0 & 0 \\ -1/R_x^{2} & 0 \end{bmatrix},
\end{align}
and $\partial T_f/\partial\alpha = \partial T_{\text{load}}/\partial\theta = 0$.

\paragraph{Chain rule.}
Apply the product rule to $T = T_1 T_f T_2 T_{\text{load}}$:
\begin{align}
\frac{\partial T}{\partial \alpha}
   &\;=\; \frac{\partial T_1}{\partial \alpha}\, T_f T_2 T_{\text{load}}
       \;+\; T_1 T_f \,\frac{\partial T_2}{\partial \alpha}\, T_{\text{load}},
\\[4pt]
\frac{\partial T}{\partial R_x}
   &\;=\; T_1 \,\frac{\partial T_f}{\partial R_x}\, T_2 T_{\text{load}}.
\end{align}
Only two of the four product-rule terms are nonzero for
$\partial/\partial\alpha$ (because $T_f$ and $T_{\text{load}}$ do
not depend on $\alpha$); only one term is nonzero for
$\partial/\partial R_x$.

\paragraph{Sending-end admittance partial.}
By the quotient rule applied to $H = C_{\text{end}}/A_{\text{end}}$:
\begin{equation}
\boxed{\;
\frac{\partial H}{\partial \theta} \;=\;
\frac{
   \dfrac{\partial C_{\text{end}}}{\partial \theta}\, A_{\text{end}}
   \;-\;
   C_{\text{end}}\, \dfrac{\partial A_{\text{end}}}{\partial \theta}
}{A_{\text{end}}^{\,2}}
\;}
\quad \theta\in\{\alpha, R_x\}.
\label{eq:dHdtheta_quotient}
\end{equation}

\paragraph{Use sites.}
The closed-form gradient (\ref{eq:dHdtheta_quotient}) is consumed by:
\begin{itemize}[leftmargin=2em]
  \item the Phase-2 swap of the optimiser's central-FD gradient
        (\texttt{inverse\_estimation/faultloc\_two\_stage\_optimiser.py},
        WP2.4 work-package);
  \item the Fisher information matrices in
        \texttt{inverse\_estimation/faultloc\_crlb\_proper.py} and
        \texttt{inverse\_estimation/faultloc\_crlb\_dualchannel.py}
        (P1.6 currently uses the central-FD partials; the swap is
        a one-line change once the closed-form module is on the
        runtime path).
\end{itemize}
The implementation is in
\texttt{inverse\_estimation/faultloc\_analytical\_gradients.py};
\texttt{tests/test\_analytical\_gradients.py} verifies a relative
error below $10^{-4}$ vs central FD at the three brief test points
$(\alpha, R_x) \in \{(0.2, 100), (0.5, 1000), (0.8, 5000)\}$, and
\texttt{matlab/tests/test\_distributed\_partials.m} cross-checks
against a MATLAB \texttt{sym/diff} derivation to $10^{-9}$.

\bigskip\hrule\bigskip
\noindent\textit{End of Appendix~A.}

\end{document}
